
% Question 14
\question
The height and the base radius of a solid metal right circular cone are $45\text{ cm}$ and $24\text{ cm}$ respectively. The circular cone is divided into a circular cone and a frustum $X$ by a plane parallel to its base. It is given that the height of $X$ is $30\text{ cm}$.
\begin{parts}
\part Express the volume of $X$ in terms of $\pi$. \rmarks{3}
\part Express the total surface area of $X$ in terms of $\pi$. \rmarks{3}
\part $X$ is melted and recast into a solid cube. Does the total surface area of the cube exceed the total surface area of $X$? Explain your answer. \rmarks{3}
\end{parts}

\begin{solution}
\begin{parts}
\part The height of the smaller circular cone removed is $45 - 30 = 15\text{ cm}$. \\
By similar figures, the base radius of the smaller circular cone is
\[
24 \times \frac{15}{45} = 8\text{ cm} \quad \text{(1M)}
\]
The volume of $X$ is
\begin{align*}
&= \frac{1}{3}\pi (24)^2(45) - \frac{1}{3}\pi (8)^2(15) \quad &\text{(1M)} \\
&= 8640\pi - 320\pi \\
&= 8320\pi\text{ cm}^3 \quad &\text{(1A)}
\end{align*}
\part The slant height of the original cone is $\sqrt{45^2 + 24^2} = 51\text{ cm}$, and the slant height of the smaller cone is $\sqrt{15^2 + 8^2} = 17\text{ cm}$. \hfill (1M) \\
The total surface area of $X$ is
\begin{align*}
&= [\pi(24)(51) - \pi(8)(17)] + \pi(24)^2 + \pi(8)^2 \quad &\text{(1M)} \\
&= 1088\pi + 576\pi + 64\pi \\
&= 1728\pi\text{ cm}^2 \quad &\text{(1A)}
\end{align*}
\part Let $s\text{ cm}$ be the side length of the solid cube.
\[
s^3 = 8320\pi \implies s = \sqrt[3]{8320\pi} \approx 29.67757 \quad \text{(1M)}
\]
The total surface area of the cube is
\[
6s^2 = 6(8320\pi)^{\frac{2}{3}} \approx 5284.45\text{ cm}^2 \quad \text{(1M)}
\]
Note that the total surface area of $X$ is $1728\pi \approx 5428.67\text{ cm}^2 > 5284.45\text{ cm}^2$. \\
Thus, the total surface area of the cube does not exceed the total surface area of $X$. \hfill (1A f.t.)
\end{parts}
\end{solution}